% !TEX root = ../../tree_notes.tex
% =====================================================================
%  Chapter 17 -- Union-Find and Minimum Spanning Trees
%  Status: OUTLINE (not written yet). See structure/chapter.txt.
% =====================================================================
\bookchapter{Union-Find and Minimum Spanning Trees}{c17}
\chapterinfo{CF 1500--2100}{Outline}

\plannedtopic{Disjoint Set Union}{9}{%
  \item A forest of parent pointers; \code{find} and \code{union}
  \item Path compression and union by size/rank
  \item Almost-$O(1)$ amortised: the inverse Ackermann function
  \item Detecting a cycle while adding edges
}

\plannedtopic{Minimum Spanning Tree}{9}{%
  \item Spanning tree of a connected graph; $n-1$ edges
  \item Kruskal with DSU, and the cut property
  \item Prim's algorithm with a priority queue
  \item MST on a tree is the tree itself
}

\plannedtopic{DSU Trees and Offline Tricks}{8}{%
  \item DSU tree (Kruskal reconstruction tree)
  \item Small-to-large merging on trees
  \item DSU on tree for subtree queries
  \item Rollback DSU and divide-and-conquer over time
}

\plannedtopic{Pruefer Sequences and Counting Trees}{9}{%
  \item The Pruefer code: a labeled tree on $n$ vertices $\leftrightarrow$ a sequence of length $n-2$
  \item Encoding and decoding in $O(n)$ / $O(n\log n)$
  \item Cayley's formula $n^{\,n-2}$ as a one-line corollary (home of the bijection; \chref{c01} states the result)
  \item Kirchhoff's matrix-tree theorem: counting spanning trees of any graph
  \item Counting trees and forests with degree constraints via multinomials
}

\plannedtopic{Spanning Tree Practice}{7}{%
  \item Second-best MST (the idea)
  \item Counting spanning trees: Cayley and Kirchhoff (matrix tree theorem)
  \item Number of connected components after deletions
}

\closechapter
